% lp2graph canonical LaTeX — assembled by corpusbuilder.repo_corpus
% Source repository: lintim/openlintim
%@ meta id=lintim__openlintim__line-planning-cost-model family=milp schema=0.1.0
%@ name :: line-planning-cost-model
%@ desc :: Cost model of line planning: choose an integer frequency for each candidate line in a given line pool so that the total operating cost is minimized while every edge of the public transportation network (PTN) carries a number of line frequencies within its lower and upper frequency bounds. Pure integer program over line frequencies.
%@ prov source :: https://gitlab.com/lintim/openlintim
%@ prov author :: lintim/openlintim
%@ index L ordered=0 cyclic=0 :: set of candidate lines in the line pool
%@ index E ordered=0 cyclic=0 :: set of edges (links) of the PTN
%@ param cost shape=L kind=vector domain=- :: operating cost of line l (used as objective coefficient on its frequency) [domain R_>=0]
%@ param f_min shape=E kind=vector domain=- :: lower frequency bound on edge e [domain Z_>=0]
%@ param f_max shape=E kind=vector domain=- :: upper frequency bound on edge e [domain Z_>=0]
%@ var f shape=L domain=integer role=primary drole=- lo=- hi=- :: frequency of line l (number of times line l is operated per period)
%@ obj sense=min name=objective combination=sum :: minimize total operating cost = sum over lines of (line cost * line frequency)
%@ con edge_lower_frequency_bound kind=linear domain=- indicator=- :: for each PTN edge, the sum of frequencies of lines passing over that edge is at least the edge lower frequency bound [capture: sum restriction not representable: 'e \\in l']
%@ con edge_upper_frequency_bound kind=linear domain=- indicator=- :: for each PTN edge, the sum of frequencies of lines passing over that edge is at most the edge upper frequency bound [capture: sum restriction not representable: 'e \\in l']
\begin{align}
  \min\quad & \sum_{l \in \mathcal{L}} \mathit{cost} \cdot f_{l} \tag{objective} \\
  & \sum_{l \in \mathcal{L}} f_{l} \ge \mathit{f\_min}_{e} \qquad \forall e \in \mathcal{E} \tag{edge\_lower\_frequency\_bound} \\
  & \sum_{l \in \mathcal{L}} f_{l} \le \mathit{f\_max}_{e} \qquad \forall e \in \mathcal{E} \tag{edge\_upper\_frequency\_bound} \\
\end{align}
